\section{SAM 系列：从实例分割到语义分割}

\subsection{SAM 1 与 SAM 2：实例级分割}

SAM（Segment Anything Model）是 Meta 推出的通用分割模型系列。SAM 1 和 SAM 2 都支持实例级别的分割，主要输入方式包括：

\begin{itemize}
    \item \textbf{Point Prompt}：在图片中指定一个点，分割该点所在的主体（如鲸鱼身体中心的点，分割出整条鲸鱼）
    \item \textbf{Bounding Box Prompt}：给定一个边界框，分割框内的主体
    \item \textbf{Automatic Mask Generation}：不需要任何输入提示，自动对整张图片进行全图实例分割
\end{itemize}

\begin{knowledgebox}{SAM 的 Point/Box 输入本身就是 Visual Prompting}
从广义上看，SAM 接受的 Point 和 Bounding Box 输入，本身就是一种 Visual Prompting——在原始图片中指定具体的视觉位置信息，引导模型进行分割。这种设计使得 SAM 可以灵活地与其他模型组合使用。
\end{knowledgebox}

SAM 1 和 SAM 2 的自动分割功能（Automatic Segmentation）不需要任何人工输入，只接收原始图片，就能对全图所有实例进行分割。这一功能是 Set-of-Mark 等 Visual Prompting 方法的基础组件。

\subsection{SAM 3：语义级可提示分割}

SAM 3 相比前代有了质的飞跃，支持语义级别的分割，且是 promptable（可提示）的：

\begin{importantbox}{SAM 3 的多模态输入}
SAM 3 支持以下多种输入方式：
\begin{itemize}
    \item \textbf{文本输入}：输入一个简单的名词短语（如 ``cat''），分割出图中所有的猫
    \item \textbf{Point/Box 输入}：继承自 SAM 1/2 的经典交互方式
    \item \textbf{示例图输入}：提供正样本和负样本图片，向正样本靠近、远离负样本
\end{itemize}
\end{importantbox}

\begin{warningbox}{SAM 3 的文本输入限制}
SAM 3 的文本输入只支持简单的名词短语（noun phrase），\textbf{不支持方位词}（如"左边的猫"、"最上面的杯子"等包含空间关系的描述）。如果需要处理包含方位描述的复杂查询，需要借助 VLM 先将复杂描述简化为名词短语，再交给 SAM 3。
\end{warningbox}

\subsection{VLM + SAM 3 Agent 组合}

由于 SAM 3 不支持方位词，可以构建一个 VLM + SAM 3 的 Agent Pipeline 来处理复杂的物体定位任务：

\begin{enumerate}
    \item \textbf{用户输入复杂查询}：如"穿蓝色马甲的最左侧的孩子"
    \item \textbf{VLM 简化查询}：将复杂描述转化为简单名词短语——"child in blue vest"
    \item \textbf{SAM 3 分割}：基于简化后的文本，分割出所有穿蓝色马甲的孩子
    \item \textbf{打上编号}：对分割结果打上数字编号
    \item \textbf{VLM 二次推理}：基于带编号的分割图，VLM 可以准确定位"最左侧的"那个孩子
\end{enumerate}

这种 Pipeline 的核心思想是：VLM 负责语言理解和推理，SAM 3 负责视觉分割，两者各司其职，协同完成单一模型难以完成的复杂定位任务。

\subsection{本章小结}

SAM 系列从 SAM 1/2 的实例分割，演进到 SAM 3 的语义级可提示分割。SAM 3 支持文本、点、框、示例图等多模态输入，但文本仅限简单名词短语。通过 VLM + SAM 3 的 Agent 组合，可以处理包含方位词的复杂定位任务。


